% Created 2026-02-04 Wed 20:55
% Intended LaTeX compiler: pdflatex
\documentclass[hyperref={pdfpagelabels=false},aspectratio=169]{beamer}
\usepackage[utf8]{inputenc}
\usepackage[T1]{fontenc}
\usepackage{graphicx}
\usepackage{longtable}
\usepackage{wrapfig}
\usepackage{rotating}
\usepackage[normalem]{ulem}
\usepackage{amsmath}
\usepackage{amssymb}
\usepackage{capt-of}
\usepackage{hyperref}
\usetheme{Madrid}
\usecolortheme{beaver}
\useinnertheme{rounded}
\useoutertheme{tree}
\author{Cristian A. Marocico}
\date{\today}
\title{Machine Learning in Python}
\subtitle{Neural Networks I: Foundations and Optimization}
\author[Marocico, Tatar]{Cristian A. Marocico, A. Emin Tatar}
\institute[CIT]{Center for Information Technology\\University of Groningen}
\RequirePackage{pgfcore}
\usepackage{amsfonts}
\usepackage{amsmath}
\usepackage{amssymb}
\usepackage{blkarray}
\usepackage{eurosym}
\usepackage[english]{babel}
\usepackage{xcolor,colortbl}
\usepackage[utf8]{inputenc}
\usepackage[OT1]{fontenc}
\usepackage{multirow}
\usepackage{listings}
\usepackage{fancyvrb}
\usepackage{tikz}
\usepackage{graphicx}
\usepackage[makeroom]{cancel}
\usepackage{bbm}
\newcommand{\infoTitle}[1]{\renewcommand{\givenTitle}{#1}}
\newcommand{\givenTitle}{Info}
\newenvironment{warning}[1][Info]{%
\infoTitle{#1}
\setbeamercolor{block title}{fg=white,bg=red!100!white}%
\setbeamercolor{block body}{bg=red!10!white}
\begin{block}{\givenTitle}}{
\end{block}}
\DeclareMathOperator*{\argmax}{argmax}
\setbeamercovered{transparent}
\usecolortheme{beaver}
\RequirePackage{pgfcore}
\setbeamercovered{transparent=1}
\mode<presentation>{
\usecolortheme{beaver}
\definecolor{rugcolor}{rgb}{0.8,0,0}
\definecolor{darkblue}{rgb}{0.13,0.29,0.53}
\definecolor{darkgreen}{rgb}{0.0,0.43,0.0}
\definecolor{darkyellow}{rgb}{0.0,0.43,0.43}
\definecolor{codegreen}{rgb}{0,0.6,0}
\definecolor{codegray}{rgb}{0.5,0.5,0.5}
\definecolor{codepurple}{rgb}{0.58,0,0.82}
\definecolor{backcolour}{rgb}{0.95,0.95,0.92}
\definecolor{Gray}{gray}{0.85}
\beamertemplatenavigationsymbolsempty
\setbeamercolor{item}{fg=rugcolor!80!black}
\setbeamercolor{title}{fg=rugcolor!80!black}
\setbeamercolor{frametitle}{fg=rugcolor!80!black, bg=black!10!white}
% Colors for 'definition' environment
\setbeamercolor*{block title}{fg=white, bg=darkblue}
\setbeamercolor*{block body}{bg=normal text.bg!85!darkblue}
% Color for the 'question' environment
\setbeamercolor*{block title question}{fg=white, bg=darkyellow}
\setbeamercolor*{block body question}{bg=normal text.bg!85!darkyellow}
\setbeamercolor*{palette tertiary}{bg=rugcolor,fg=white}
%HEADER WITH HIGHLIGHTED SECTION NAMES (optional)
\useheadtemplate{%
\vbox{%
%			\vskip1.2pt
%			\pgfuseimage{logo}
%			\vskip1.2pt
\tinycolouredline{rugcolor}{
\color{white}{
% comment the following line if you don't want the section names lines
%					to appear on top
\insertsectionnavigationhorizontal{\paperwidth}{}{\hskip0pt
plus1filll}
%\pgfuseimage{logored}
}
}
%    \tinycolouredline{rugcolor}
{\color{white}{
%% \insertsectionnavigationhorizontal{\paperwidth}{}{
%                \hskip0pt \hfill}
}}
}
}
%FOOTER WITH AUTHOR NAME(S), PAPER TITLE (ABBREVIATED IF SPECIFIED BY \title),
% AND PAGE COUNTER (optional)
%	\usefoottemplate{%
%		\vbox{%
%			\tinycolouredline{rugcolor}{
%				\color{white}{
%					{\insertshortauthor} \hfill \insertshortsubtitle \hfill
%					%\insertdate \hfill%
%					\textsc{\insertframenumber/\inserttotalframenumber}
%         		}
%         	}
%		}
%	}
}
\newtheorem*{props}{Properties}
\newtheorem*{prop}{Property}
\newtheorem*{notation}{Notation}
\newtheorem*{terminology}{Terminology}
\newcolumntype{a}{>{\columncolor{Gray}}c}
\DeclareMathOperator*{\argmin}{argmin}
\newcommand{\Var}{\mathbb{V}\mathrm{ar}}
\newcommand{\Corr}{\mathbb{C}\mathrm{orr}}
\newcommand{\Cov}{\mathbb{C}\mathrm{ov}}
\newcommand{\Expt}{\mathbb{E}}
\newcommand{\NorDist}{\mathcal{N}}
\newcommand{\ExpDist}{\mathcal{E}\mathrm{xp}}
\newcommand{\GammaDist}{\mathcal{G}\mathrm{amma}}
\newcommand{\BetaDist}{\mathcal{B}\mathrm{eta}}
\newcounter{listCounter}
\newenvironment{question}{%
\setbeamercolor{block title}{bg=orange!70!white,fg=white}
\setbeamercolor{block body}{bg=yellow!10!white}
\begin{block}{Question}
}{%
\end{block}
}
\lstdefinestyle{mystyle}{
language=R,
backgroundcolor=\color{backcolour},
commentstyle=\color{codegreen},
keywordstyle=\color{darkblue}\bfseries,
numberstyle=\tiny\color{codegray},
stringstyle=\color{codepurple},
basicstyle=\ttfamily,
breakatwhitespace=true,
breaklines=true,
%	captionpos=none,
columns=fixed,
keepspaces=true,
numbers=none,
numbersep=5pt,
showspaces=false,
showstringspaces=false,
showtabs=false,
tabsize=2,
basewidth=1.5em,
escapeinside={<@}{@>}
}
\setbeamertemplate{caption}{\raggedright\insertcaption\par}
\setbeamertemplate{blocks}[rounded][shadow=true]
\renewenvironment{definition}[1][Definition]{%
\setbeamercolor{block title}{fg=white,bg=darkblue}
\setbeamercolor{block body}{fg=black,bg=normal text.bg!85!darkblue}
\begin{block}{#1\hfill \footnotesize{Definition}}
\vspace*{-5pt}
}{%
\end{block}
}
\renewenvironment{example}[1][Example]{%
% Color for the 'example' environment
\setbeamercolor*{block title}{fg=white, bg=darkgreen}
\setbeamercolor*{block body}{bg=normal text.bg!85!darkgreen}
\begin{block}{#1\hfill \footnotesize{Example}}
\vspace*{-5pt}
}{%
\end{block}
}
\renewenvironment{theorem}[1][Theorem]{%
\setbeamercolor*{block title}{fg=white, bg=darkyellow}
\setbeamercolor*{block body}{bg=normal text.bg!85!darkyellow}
\begin{block}{#1\hfill \footnotesize{Theorem}}
\vspace*{-5pt}
}{%
\end{block}
}
\date[Feb 17\textsuperscript{th} 2026]{Tuesday, February 17\textsuperscript{th} 2026}
\usepackage{mathtools}
\newcommand{\intsum}{\mathop{\mathrlap{\raisebox{0.1ex}{\hspace{0.2em}$\textstyle\sum$}}\int}\limits}
\setbeamercovered{transparent=0}
\usepackage[timeinterval=60]{tdclock}
\hypersetup{
 pdfauthor={Cristian A. Marocico},
 pdftitle={Machine Learning in Python},
 pdfkeywords={},
 pdfsubject={},
 pdfcreator={Emacs 30.2 (Org mode 9.7.11)}, 
 pdflang={English}}
\begin{document}

\maketitle
\begin{frame}{Outline}
\setcounter{tocdepth}{1}
\tableofcontents
\end{frame}

\section{Introduction to Artificial Neural Networks}
\label{sec:org04297b1}
\begin{frame}[label={sec:orgcd8d6cb}]{What is an Artificial Neural Network?}
\begin{definition}[Definition]\label{sec:org425637c}
\pause
An \alert{Artificial Neural Network (ANN)} is a computational model composed of layers of interconnected nodes ("neurons") that process information by transmitting signals to one another.
\pause
\end{definition}
\begin{example}[Key Properties]\label{sec:org264a45f}
\begin{itemize}[<+->]
\item Foundation of \alert{Deep Learning}
\item Can learn complex, non-linear relationships
\item Powers image recognition, voice assistants, medical diagnosis, and more
\end{itemize}
\end{example}
\end{frame}
\begin{frame}[label={sec:orgc94818c}]{Key Components of a Neural Network}
\begin{block}{The Building Blocks}
\pause
\begin{enumerate}[<+->]
\item \alert{Neurons (Nodes)}: Basic processing units that receive inputs, perform a calculation, and produce an output
\item \alert{Weights}: Connections between neurons that determine influence; learning = adjusting weights
\item \alert{Layers}:
\begin{itemize}
\item \alert{Input Layer}: Receives raw data
\item \alert{Hidden Layers}: Feature extraction (the "magic")
\item \alert{Output Layer}: Final prediction
\end{itemize}
\item \alert{Activation Function}: Introduces non-linearity, allowing the network to learn curves
\end{enumerate}
\end{block}
\end{frame}
\begin{frame}[label={sec:org30143e6}]{Biological Inspiration}
\begin{block}{From Biology to Mathematics}
\pause
ANNs are loosely inspired by biological neural networks in the brain:
\pause
\begin{center}
\begin{tabular}{lll}
Biological Neuron & Artificial Neuron & Role\\
\hline
\alert{Dendrites} & Inputs (\(x\)) & Receives signals\\
\alert{Synapse} & Weights (\(w\)) & Controls signal strength\\
\alert{Cell Body} & Summation (\(\sum\)) & Aggregates signals\\
\alert{Axon} & Activation & Transmits output\\
\end{tabular}
\end{center}
\end{block}
\end{frame}
\begin{frame}[label={sec:org984a1bc}]{The Perceptron (1958)}
\begin{definition}[The Simplest Neural Network]\label{sec:orgee6bc15}
\pause
The \alert{Perceptron} is a single-layer network that takes inputs, weighs them, and makes a binary decision.
\pause
\end{definition}
\begin{block}{Mathematical Formulation}
\begin{enumerate}
\item \alert{Weighted Sum}: \(z = w_1 x_1 + w_2 x_2 + \dots + b\)
\item \alert{Activation}:
$$\text{Output} = \begin{cases} 1 & \text{if } z > 0 \\ 0 & \text{if } z \le 0 \end{cases}$$
\end{enumerate}
\pause
This can solve \alert{linearly separable} problems (drawing a straight line between classes).
\end{block}
\end{frame}
\begin{frame}[label={sec:orgd231bbb}]{The XOR Problem}
\begin{block}{The Fatal Flaw of Single Perceptrons}
\pause
\begin{center}
\begin{tabular}{rrrr}
\(x_1\) & \(x_2\) & AND & XOR\\
\hline
0 & 0 & 0 & 0\\
0 & 1 & 0 & 1\\
1 & 0 & 0 & 1\\
1 & 1 & 1 & 0\\
\end{tabular}
\end{center}
\pause
\end{block}
\begin{example}[The Limitation]\label{sec:org5c11fa3}
\begin{itemize}[<+->]
\item \alert{AND Gate}: Linearly separable (one line can separate classes)
\item \alert{XOR Gate}: NOT linearly separable (needs a curve or multiple lines)
\item This limitation froze AI research until \alert{Multi-Layer Perceptrons (MLPs)} emerged
\end{itemize}
\end{example}
\end{frame}
\section{The Multilayer Perceptron (MLP)}
\label{sec:orgc289329}
\begin{frame}[label={sec:org0461d23}]{From One to Many}
\begin{definition}[Stacking Neurons]\label{sec:orgcb28804}
\pause
A \alert{Multilayer Perceptron (MLP)} or \alert{Feedforward Neural Network} connects many neurons in layers to solve complex, non-linear problems.
\pause
\end{definition}
\begin{block}{Architecture}
\begin{enumerate}[<+->]
\item \alert{Input Layer}: Receives raw feature vector \(X\) (no computation)
\item \alert{Hidden Layers}: Perform computation and extract features; more layers = "deeper" network
\item \alert{Output Layer}: Produces final prediction (e.g., probability of class)
\end{enumerate}
\end{block}
\end{frame}
\begin{frame}[label={sec:org3583b53}]{Activation Functions}
\begin{block}{Why Activation Functions?}
\pause
Without activation functions, an MLP is just a giant Linear Regression (linear + linear = linear).
\alert{Activation functions} introduce \alert{non-linearity}, allowing networks to learn curves.
\pause
\end{block}
\begin{example}[Common Activation Functions]\label{sec:orga684aa9}
\begin{center}
\begin{tabular}{llll}
Function & Formula & Range & Usage\\
\hline
\alert{Sigmoid} & \(\frac{1}{1+e^{-z}}\) & \((0, 1)\) & Output (Binary Class.)\\
\alert{Tanh} & \(\frac{e^z - e^{-z}}{e^z + e^{-z}}\) & \((-1, 1)\) & Hidden layers (Legacy)\\
\alert{ReLU} & \(\max(0, z)\) & \([0, \infty)\) & \alert{Standard for Hidden}\\
\alert{Softmax} & \(\frac{e^{z_i}}{\sum e^{z_j}}\) & \((0, 1)\) & Output (Multiclass)\\
\end{tabular}
\end{center}
\end{example}
\end{frame}
\begin{frame}[label={sec:orgf3acf1b},fragile]{Forward Propagation}
 \begin{block}{How Data Flows Through the Network}
\pause
For a single layer \(l\):
\pause
\begin{enumerate}[<+->]
\item \alert{Linear Step (\(Z\))}: Weighted sum of inputs from previous layer
$$Z^{[l]} = W^{[l]} \cdot A^{[l-1]} + b^{[l]}$$
\item \alert{Activation Step (\(A\))}: Apply activation function
$$A^{[l]} = g(Z^{[l]})$$
\end{enumerate}
\pause
This process repeats until the output layer. The magic: \texttt{np.dot} handles all connections automatically.
\end{block}
\end{frame}
\begin{frame}[label={sec:org21c385f}]{Network Capacity: Width vs Depth}
\begin{block}{Architectural Choices}
\pause
\begin{center}
\begin{tabular}{lll}
Architecture & Structure & Characteristics\\
\hline
\alert{Wide} & 1 layer, 100 neurons & Simple patterns, fast\\
\alert{Deep} & 5 layers, 10 neurons & Hierarchical features, complex\\
\end{tabular}
\end{center}
\pause
\end{block}
\begin{example}[Trade-offs]\label{sec:orgdd4a53d}
\begin{itemize}[<+->]
\item Deeper networks learn hierarchical representations (edges \(\rightarrow\) parts \(\rightarrow\) objects)
\item Wider networks can approximate functions but may need more parameters
\item In practice: start simple, add complexity if needed
\end{itemize}
\end{example}
\end{frame}
\section{Training Neural Networks}
\label{sec:orge8d9a82}
\begin{frame}[label={sec:org2891de2}]{The Goal of Training}
\begin{block}{From Random to Learned}
\pause
A neural network starts with \alert{random weights}. It knows nothing.
\alert{Training} adjusts these weights to minimize prediction error.
\pause
\end{block}
\begin{example}[Three Components]\label{sec:org63cf105}
\begin{enumerate}[<+->]
\item \alert{Loss Function}: How we measure error
\item \alert{Optimizer}: How we update weights
\item \alert{Backpropagation}: How we calculate the updates
\end{enumerate}
\end{example}
\end{frame}
\begin{frame}[label={sec:org30387d1}]{Loss Functions}
\begin{definition}[The Scorecard]\label{sec:orgac2f699}
\pause
The \alert{Loss Function} (Cost Function) calculates a single number representing "how bad" the model is. Goal: drive it to zero.
\pause
\end{definition}
\begin{block}{Common Loss Functions}
\begin{center}
\begin{tabular}{lll}
Problem Type & Loss Function & Formula\\
\hline
\alert{Regression} & Mean Squared Error & \(\frac{1}{n}\sum(y - \hat{y})^2\)\\
\alert{Classification} & Cross-Entropy & \(-\sum y \log(\hat{y})\)\\
\end{tabular}
\end{center}
\pause
\begin{itemize}
\item MSE: Penalizes large errors heavily (prices, temperatures)
\item Cross-Entropy: Penalizes confident wrong answers (probabilities)
\end{itemize}
\end{block}
\end{frame}
\begin{frame}[label={sec:org30588bb}]{Stochastic Gradient Descent (SGD)}
\begin{block}{Optimization Strategy}
\pause
\alert{Gradient Descent} minimizes loss by following the slope downhill.
\alert{Batch GD} calculates gradient for entire dataset — slow and memory-intensive.
\pause
\end{block}
\begin{definition}[Stochastic Gradient Descent]\label{sec:org4ec2235}
\alert{SGD} updates weights after seeing just \alert{one example} (or a small "mini-batch").
\pause
\begin{itemize}[<+->]
\item \alert{Pros}: Faster, noise helps escape local minima
\item \alert{Cons}: Zigzaggy path to convergence
\end{itemize}
\end{definition}
\end{frame}
\begin{frame}[label={sec:org45e5301}]{Backpropagation}
\begin{block}{The Chain Rule in Action}
\pause
How do we know which weight in Layer 1 caused the error in Layer 5?
We use the \alert{Chain Rule} of calculus.
\pause
\end{block}
\begin{example}[The Process]\label{sec:org57f346c}
\begin{enumerate}[<+->]
\item \alert{Forward Pass}: Data flows through network, loss is calculated
\item \alert{Backward Pass}: Gradient propagates backwards layer by layer
\begin{itemize}
\item "Layer 4, you contributed X to the error. Adjust your weights."
\item "Layer 3, based on Layer 4's error, you contributed Y\ldots{}"
\end{itemize}
\end{enumerate}
\end{example}
\end{frame}
\begin{frame}[label={sec:org0d588de}]{The Vanishing Gradient Problem}
\begin{block}{Why Deep Networks Failed (Initially)}
\pause
In deep networks, gradients are multiplied repeatedly during backpropagation.
\alert{Sigmoid} has max derivative of 0.25:
$$0.25 \times 0.25 \times 0.25 \times \dots \approx 0$$
\pause
\end{block}
\begin{example}[The Consequence]\label{sec:org2c10375}
\begin{itemize}[<+->]
\item Gradient "vanishes" before reaching early layers
\item Early layers never learn; deep network fails
\item \alert{Solution}: Use \alert{ReLU} (derivative = 1 for positive inputs)
\end{itemize}
\end{example}
\end{frame}
\begin{frame}[label={sec:org95e930c}]{The Learning Rate}
\begin{block}{The Most Critical Hyperparameter}
\pause
The \alert{Learning Rate} (\(\alpha\)) controls step size during gradient descent.
\pause
\end{block}
\begin{example}[The Goldilocks Zone]\label{sec:org351e223}
\begin{center}
\begin{tabular}{ll}
Learning Rate & Behavior\\
\hline
\alert{Too Small} & Loss decreases very slowly (0.0001)\\
\alert{Too Large} & Loss oscillates or explodes (10.0)\\
\alert{Just Right} & Smooth, rapid descent to minimum (0.01-0.1)\\
\end{tabular}
\end{center}
\end{example}
\end{frame}
\section{Optimizing Neural Networks}
\label{sec:org07e48ca}
\begin{frame}[label={sec:org656381a}]{Weight Initialization}
\begin{block}{Starting Right}
\pause
\begin{center}
\begin{tabular}{ll}
Initialization & Problem\\
\hline
All Zeros & Symmetry: every neuron learns same thing\\
Too Large & Gradients explode\\
Too Small & Gradients vanish\\
\end{tabular}
\end{center}
\pause
\end{block}
\begin{example}[Strategies]\label{sec:org806a4f0}
\begin{itemize}[<+->]
\item \alert{He Initialization}: Best for ReLU networks
\item \alert{Xavier (Glorot)}: Best for Sigmoid/Tanh networks
\item Both keep variance of activations constant across layers
\end{itemize}
\end{example}
\end{frame}
\begin{frame}[label={sec:org700831d}]{Advanced Optimizers}
\begin{block}{Smarter Descent}
\pause
Standard SGD is like walking down a mountain blindfolded. Advanced optimizers add intelligence.
\pause
\begin{center}
\begin{tabular}{lll}
Optimizer & Mechanism & Analogy\\
\hline
\alert{SGD + Momentum} & Accumulates past gradients & Heavy ball rolling downhill\\
\alert{RMSProp} & Adapts learning rate per parameter & Slows on steep, speeds on flat\\
\alert{Adam} & Momentum + RMSProp combined & Smart ball, adapts speed \& direction\\
\end{tabular}
\end{center}
\pause
\alert{Adam} is the default choice for most tasks.
\end{block}
\end{frame}
\begin{frame}[label={sec:org94ff7b3}]{Regularization: Dropout}
\begin{definition}[Fighting Overfitting]\label{sec:org05e9325}
\pause
\alert{Dropout} randomly "kills" (sets to zero) a percentage of neurons during training.
\pause
\end{definition}
\begin{example}[Why It Works]\label{sec:org3ac43f0}
\begin{itemize}[<+->]
\item Prevents neurons from co-adapting
\item Forces network to be redundant and robust
\item Each training pass uses a different "thinned" network
\item At test time, all neurons are used (with scaled weights)
\end{itemize}
\end{example}
\end{frame}
\begin{frame}[label={sec:org7dc11f4}]{Regularization: Batch Normalization}
\begin{definition}[Stabilizing Training]\label{sec:org042a064}
\pause
\alert{Batch Normalization} normalizes inputs of each layer to have mean 0 and variance 1 during training.
\pause
\end{definition}
\begin{example}[Benefits]\label{sec:org3248e9f}
\begin{itemize}[<+->]
\item Stabilizes training (allows higher learning rates)
\item Reduces sensitivity to initialization
\item Acts as a form of regularization
\item Enables training of much deeper networks
\end{itemize}
\end{example}
\end{frame}
\begin{frame}[label={sec:org97fbc0f}]{Early Stopping}
\begin{block}{The Simplest Regularization}
\pause
Monitor the \alert{Validation Loss}. If it stops improving (or starts getting worse), stop training immediately.
\pause
\end{block}
\begin{example}[How It Works]\label{sec:org51198f3}
\begin{itemize}[<+->]
\item Training loss keeps decreasing (memorization)
\item Validation loss starts increasing (overfitting)
\item Early stopping saves the model at the sweet spot
\item No hyperparameter tuning needed (just patience parameter)
\end{itemize}
\end{example}
\end{frame}
\section{Building Networks in PyTorch}
\label{sec:org0a8b05a}
\begin{frame}[label={sec:org0b7b644}]{Why PyTorch?}
\begin{block}{Beyond Scikit-learn}
\pause
Scikit-learn is great for standard ML, but Deep Learning needs:
\pause
\begin{enumerate}[<+->]
\item \alert{GPU Acceleration}: Matrix multiplication is slow on CPUs
\item \alert{Automatic Differentiation}: Backpropagation is hard to code by hand
\end{enumerate}
\pause
\end{block}
\begin{example}[PyTorch Advantages]\label{sec:org6f1abea}
\begin{itemize}
\item Readable, "Pythonic" style
\item Research and education standard
\item Dynamic computation graphs (easy debugging)
\end{itemize}
\end{example}
\end{frame}
\begin{frame}[label={sec:org816bc5d},fragile]{Tensors: The Foundation}
 \begin{definition}[NumPy Arrays on Steroids]\label{sec:org2637cd9}
\pause
A \alert{Tensor} is like a NumPy array that:
\begin{itemize}
\item Can live on a GPU
\item Remembers how it was created (for backprop)
\end{itemize}
\pause
\end{definition}
\begin{example}[Key Operations]\label{sec:orgaa85b47}
\begin{itemize}[<+->]
\item \texttt{torch.tensor([1, 2, 3])} — Create from list
\item \texttt{torch.ones(2, 3)} — Shape (2, 3) of ones
\item \texttt{x.to("cuda")} — Move to GPU
\item Math: \texttt{x * 2 + 5} — Same as NumPy
\end{itemize}
\end{example}
\end{frame}
\begin{frame}[label={sec:orgf18e81b},fragile]{Defining a Network (\texttt{nn.Module})}
 \begin{block}{The Blueprint}
\pause
In PyTorch, a network is a Python class inheriting from \texttt{nn.Module}:
\pause
\end{block}
\begin{example}[Two Required Methods]\label{sec:orgf40c96f}
\begin{enumerate}[<+->]
\item \texttt{\_\_init\_\_}: Define the layers
\begin{itemize}
\item \texttt{self.layer1 = nn.Linear(input\_dim, hidden\_dim)}
\item \texttt{self.activation = nn.ReLU()}
\end{itemize}
\item \texttt{forward}: Define data flow
\begin{itemize}
\item \texttt{x = self.layer1(x)}
\item \texttt{x = self.activation(x)}
\item \texttt{return x}
\end{itemize}
\end{enumerate}
\end{example}
\end{frame}
\begin{frame}[label={sec:org88e8805},fragile]{The Training Loop}
 \begin{block}{The 5-Step Boilerplate}
\pause
Unlike Scikit-learn's \texttt{model.fit()}, PyTorch requires manual control:
\pause
\begin{enumerate}[<+->]
\item \alert{Forward Pass}: \texttt{y\_pred = model(X)}
\item \alert{Compute Loss}: \texttt{loss = criterion(y\_pred, y)}
\item \alert{Zero Gradients}: \texttt{optimizer.zero\_grad()}
\item \alert{Backward Pass}: \texttt{loss.backward()}
\item \alert{Update Weights}: \texttt{optimizer.step()}
\end{enumerate}
\pause
Wrap in a \texttt{for epoch in range(epochs):} loop.
\end{block}
\end{frame}
\begin{frame}[label={sec:org7821349},fragile]{Loss Functions and Optimizers}
 \begin{block}{PyTorch Components}
\pause
\end{block}
\begin{example}[Loss Functions]\label{sec:orgba42141}
\begin{center}
\begin{tabular}{ll}
Task & Loss Class\\
\hline
Binary Classification & \texttt{nn.BCELoss()}\\
Multiclass & \texttt{nn.CrossEntropyLoss()}\\
Regression & \texttt{nn.MSELoss()}\\
\end{tabular}
\end{center}
\pause
\end{example}
\begin{block}{Optimizers}
\begin{verbatim}
optimizer = torch.optim.Adam(model.parameters(), lr=0.01)
\end{verbatim}
\end{block}
\end{frame}
\begin{frame}[label={sec:orgf5619b3},fragile]{Evaluation Mode}
 \begin{block}{No Gradients Needed}
\pause
When evaluating, we don't need to calculate gradients. Use \texttt{torch.no\_grad()} to save memory:
\pause
\begin{verbatim}
with torch.no_grad():
    test_preds = model(X_test)
    predicted = (test_preds > 0.5).float()
    accuracy = (predicted == y_test).float().mean()
\end{verbatim}
\end{block}
\end{frame}
\section{Summary}
\label{sec:orgcc87042}
\begin{frame}[label={sec:org57b0d05}]{Key Takeaways}
\begin{block}{What We Learned}
\pause
\begin{enumerate}[<+->]
\item \alert{ANNs}: Layers of neurons with weights and activation functions
\item \alert{Perceptron}: Single neuron, limited to linear problems
\item \alert{MLP}: Stacked neurons solve non-linear problems (XOR)
\item \alert{Training}: Forward pass \(\rightarrow\) Loss \(\rightarrow\) Backprop \(\rightarrow\) Update
\item \alert{Vanishing Gradient}: Sigmoid fails deep; ReLU fixes it
\item \alert{Optimization}: Adam is the default; Early Stopping prevents overfitting
\item \alert{PyTorch}: Tensors + nn.Module + Training Loop
\end{enumerate}
\end{block}
\end{frame}
\begin{frame}[label={sec:orga0cb69e}]{Next Section Preview}
\begin{block}{Neural Networks II: Deep Learning \& Transformers}
\pause
\begin{itemize}[<+->]
\item \alert{Convolutional Neural Networks (CNNs)}: Image processing with filters
\item \alert{From RNNs to Transformers}: Sequence modeling revolution
\item \alert{Large Language Models}: GPT, BERT, and modern NLP
\item \alert{Generative AI}: Text generation, RAG, Diffusion models
\end{itemize}
\end{block}
\end{frame}
\end{document}
